\documentclass[a4paper,11pt]{jsarticle}

% ===== パッケージ =====
\usepackage[top=25mm, bottom=25mm, left=25mm, right=25mm]{geometry}
\usepackage{amsmath, amssymb}
\usepackage{graphicx}
\usepackage{url}
\usepackage{cite}
\usepackage{enumitem}

% ===== タイトル情報 =====
\title{レポートタイトル}
\author{氏名 \\ 東京理科大学 工学部 経営工学科 \\ 秦野研究室}
\date{\today}

\begin{document}

\maketitle

% ===== 概要 =====
\begin{abstract}
ここに概要（アブストラクト）を記述する。
200〜300字程度で研究の目的・手法・結果を簡潔にまとめる。
\end{abstract}

% ===== 1. はじめに =====
\section{はじめに}
研究の背景、目的、本レポートの構成について述べる。

大規模言語モデル（LLM）は自然言語処理の分野において
大きな進歩を遂げている\cite{sample2024}。

% ===== 2. 関連研究 =====
\section{関連研究}
関連する先行研究について述べる。

\subsection{マルチエージェントシステム}
マルチエージェントシステム（MAS）の概要を述べる。

\subsection{トークン効率化}
トークン効率化に関する先行研究を整理する。

% ===== 3. 提案手法 =====
\section{提案手法}
提案する手法について詳細に述べる。

数式の例：
\begin{equation}
  L = \sum_{i=1}^{N} \alpha_i \cdot f(x_i) + \lambda \cdot R(\theta)
\end{equation}

% ===== 4. 実験 =====
\section{実験}

\subsection{実験設定}
実験の条件、データセット、評価指標について述べる。

表の例：
\begin{table}[h]
  \centering
  \caption{実験結果の比較}
  \begin{tabular}{lcc}
    \hline
    手法 & 精度 (\%) & トークン削減率 (\%) \\
    \hline
    ベースライン & 75.3 & --- \\
    提案手法 & 78.1 & 42.5 \\
    \hline
  \end{tabular}
\end{table}

\subsection{結果と考察}
実験結果について分析・考察する。

% ===== 5. おわりに =====
\section{おわりに}
本研究のまとめと今後の課題について述べる。

% ===== 参考文献 =====
\begin{thebibliography}{9}
  \bibitem{sample2024}
    著者名,
    ``論文タイトル,''
    ジャーナル名,
    vol.~X, no.~Y, pp.~1--10,
    2024.
\end{thebibliography}

\end{document}
